\subsection*{同一个词为什么需要不同上下文表示}
想象“银行”出现在取款句子与河岸翻译任务中。初始 token 表示本身不足以决定当前语境，需要从周围位置读取信息。注意力为每个位置建立一份按上下文改变的加权汇总，而不是给所有位置固定同一平均值。

查询负责提出匹配需求，键负责供匹配，值负责携带真正传递的内容。查询与键点积得到分数，再通过归一化产生权重。键与值分开，使“怎样找到信息”和“取回什么信息”能使用不同变换。

\subsection*{用两个值理解汇总，再放回矩阵}
只有两个标量值二与十时，权重 $(0.75,0.25)$ 给出四；把第二权重增大，输出向十移动。非负且和为一的权重保证汇总位于两值之间。交互实验会同时展示分数、温度、权重和输出，便于区分这四个对象。

真实单头中，每个查询都对多个键计算一行分数，所以出现 $QK^T$。除以 $\sqrt{d_k}$ 是缩放分数的设计；因果掩码让当前位置不能偷看未来 token。对每行独立归一化，再乘 $V$，便批量完成所有查询的汇总。

凸包解释只约束这一步的值平均，后续输出投影、多个头和残差连接都会改变最终位置。查询“寻找某类词”的比喻有助于理解可能行为，但某个训练模型究竟学了什么，仍须查看参数和实验，不能从命名直接断言。
